% FAC-EN-U0038
% French source: fac_body.tex lines 1377--1383.
% Authority: PDF physical pp. 37--38; printed pp. 232--233 (no. 40 complete).

\medskip
\phantomsection\label{fac:II:2:no:40}
\textbf{40. Sheaves of fractional ideals.}
\phantomsection\label{fac:II:2:no:40:def:fractional-ideal-sheaf}
Let $V$ be an irreducible algebraic variety, and let $K(V)$ be the constant sheaf of rational functions on $V$ (cf. \hyperref[fac:II:1:no:36]{no.~36}). The sheaf $K(V)$ is algebraic, but it is not coherent if $\dim V>0$. An algebraic subsheaf $\mathcal{F}$ of $K(V)$ may be called a ``sheaf of fractional ideals,'' since each $\mathcal{F}_x$ is a fractional ideal of $\mathcal{O}_{x,V}$.

\phantomsection\label{fac:II:2:no:40:prop:4}
\textsc{Proposition 4.} \textit{An algebraic subsheaf $\mathcal{F}$ of $K(V)$ is coherent if and only if it is of finite type.}

The necessity is trivial. To prove the sufficiency, it is enough to show that $K(V)$ satisfies condition~\hyperref[fac:I:2:no:12:def:2:item:b]{(b)} of \hyperref[fac:I:2:no:12:def:2]{Definition~2 in no.~12}; in other words, if $f_1,\ldots,f_p$ are rational functions, then the sheaf $\mathcal{R}(f_1,\ldots,f_p)$ is of finite type. If $x$ is a point of $V$, we can find functions $g_i$ and $h$ such that $f_i=g_i/h$, with the $g_i$ and $h$ regular on a neighborhood $U$ of $x$ and $h$ nowhere zero on $U$. The sheaf $\mathcal{R}(f_1,\ldots,f_p)$ is then equal to $\mathcal{R}(g_1,\ldots,g_p)$, which is of finite type because $\mathcal{O}_V$ is a coherent sheaf of rings.
